Este trabalho apresenta um fluxo demonstrativo para validação automatizada de documentos acadêmicos produzidos em LaTeX. O objetivo é mostrar, em um TCC curto e estruturalmente realista, como separar a preparação do conteúdo, a compilação do projeto, a execução de verificações automáticas e a revisão visual do arquivo final. O procedimento considera propriedades objetivas que podem ser observadas no PDF, como tamanho da página, incorporação de fontes, conformidade com o alvo PDF/A adotado e presença de marcadores estruturais. Para ilustrar o registro dos resultados, utiliza-se uma execução sintética com quatro verificações e uma métrica simples de taxa de aprovação. Os dados apresentados têm finalidade exclusivamente didática e não representam um experimento externo. O trabalho também demonstra a organização prática do projeto: metadados concentrados no arquivo principal, capítulos separados, bibliografia estruturada e uso contextual de figura, tabela, equação e código-fonte. Os resultados reforçam que a automação pode antecipar problemas repetitivos e tornar a verificação mais reprodutível, mas não substitui a avaliação do conteúdo científico, da metodologia, das fontes e das exigências específicas do curso ou programa.

\ufcSummaryKeywords{LaTeX; trabalhos acadêmicos; validação; automação; normalização.}
